\newpage
\phantomsection
\addcontentsline{toc}{section}{References \& Academic Sources}
\section*{References \& Academic Sources}

The educational material, grammatical rules, and narrative anecdotes provided in this guide are deeply rooted in rigorous, authoritative scholarship.

\subsection*{The Valmiki Ramayana}
All character-building anecdotes and contextual stories are strictly sourced from the first six Kandas of the \textbf{Valmiki Ramayana}, intentionally avoiding contested texts or later additions (such as the Uttara Kanda).
\begin{itemize}
    \item \href{https://www.valmiki.iitk.ac.in}{\textbf{IIT Kanpur Valmiki Ramayana Portal}}: Our primary digital repository for exact Shloka verification and English translations.
\end{itemize}

\subsection*{The Ramcharitmanas \& Awadhi Text}
The core Devanagari text and orthography follow the traditional, universally accepted manuscripts.
\begin{itemize}
    \item \href{https://gitapress.org/result?search=Hanuman%20Chalisa&category=}{\textbf{Gita Press, Gorakhpur}}: The gold standard for the source Devanagari and Awadhi text of the Hanuman Chalisa. \href{https://dn710202.ca.archive.org/0/items/xqZf_shri-hanuman-chalisha-gita-press/Shri%20Hanuman%20Chalisha%20-%20Gita%20Press.pdf}{\textbf{PDF from ArXiv.org}}.


    \item \href{https://www.ramcharitmanas.iitk.ac.in/content/tulsidas}{\textbf{IIT Kanpur Ramcharitmanas Portal}}: Used for cross-referencing Tulsidas's broader poetic patterns and lexical choices.
\end{itemize}

\subsection*{Awadhi Grammar \& Linguistics}
The phonetic shift explanations (like the \(v \rightarrow b\) shift) and grammatical notes found in the ''Gotcha'' boxes are derived from foundational academic research in Indo-Aryan linguistics.
\begin{itemize}
    \item \textbf{Baburam Saksena, \href{https://archive.org/details/in.ernet.dli.2015.238311/mode/2up}{\textit{Evolution of Awadhi}}} (1937): The seminal academic work detailing the historical development, phonology, and morphology of the Awadhi language.
    \item \textbf{Rupert Snell, \href{https://ia600500.us.archive.org/16/items/Pushtimarg/BrajBhashaReaderenglish.pdf}{\textit{The Hindi Classical Tradition: A Braj Bhas\={a} Reader}}} (1991): The foundational text for understanding the metrical pulse (Matra/Beats) and the treatment of syllable weights in medieval North Indian poetry.
    \item \textbf{George A. Grierson, \href{https://dsal.uchicago.edu/books/lsi/lsi.php?volume=6&pages=286}{\textit{Linguistic Survey of India}}} (Vol VI): The historical baseline for dialectal mapping and grammatical structures of Eastern Hindi and Awadhi.
\end{itemize}
